\chapter{2009年辽宁卷（文）}

\section{单选题}

\begin{problem}
已知集合 \(M = \{x \mid -3 < x \leqslant 5\}\)，\(N = \{x \mid x < -5\text{ 或 }x > 5\}\)，则 \(M \cup N =\)（\quad）

\choices
    {\(\{x\mid x < -5\text{ 或 }x > -3\}\)}
    {\(\{x\mid -5 < x < 5\}\)}
    {\(\{x\mid -3 < x < 5\}\)}
    {\(\{x\mid x < -3\text{ 或 }x > 5\}\)}
\end{problem}

\begin{answer}
A
\end{answer}

\begin{problem}
已知复数 \(z = 1 - 2\i\)，那么 \(\frac{1}{z} =\)（\quad）

\choices
    {\(\frac{\sqrt{5}}{5} + \frac{2\sqrt{5}}{5}\i\)}
    {\(\frac{\sqrt{5}}{5} -\frac{2\sqrt{5}}{5}\i\)}
    {\(\frac{1}{5} +\frac{2}{5}\i\)}
    {\(\frac{1}{5} -\frac{2}{5}\i\)}
\end{problem}

\begin{answer}
C
\end{answer}

\begin{problem}
\(\{a_{n}\}\) 为等差数列，且 \(a_7 - 2a_4 = -1\)，\(a_3 = 0\)，则公差 \(d =\) （\quad）

\choices
    {\(-2\)}
    {\(-\frac{1}{2}\)}
    {\(\frac{1}{2}\)}
    {\(2\)}
\end{problem}

\begin{answer}
B
\end{answer}

\begin{problem}
平面向量 \(\bs{a}\) 与 \(\bs{b}\) 的夹角为 \( 60^{\circ} \)，\( \bs{a} = (2, 0) \)，\( |\bs{b}| = 1 \)，则 \( |\bs{a} + 2\bs{b}| = \)（\quad）

\choices
    {\(\sqrt{3}\)}
    {\(2 \sqrt{3}\)}
    {\(4\)}
    {\(12\)}
\end{problem}

\begin{answer}
B
\end{answer}

\begin{problem}
如果把地球看成一个球体，则地球上北纬 \(60^{\circ}\) 纬线和赤道线长的比值为（\quad）

\choices
    {\(0.8\)}
    {\(0.75\)}
    {\(0.5\)}
    {\(0.25\)}
\end{problem}

\begin{answer}
C
\end{answer}

\begin{problem}
已知函数 \( f(x) \) 满足：当 \( x \geqslant 4 \) 时，\( f(x) = \left(\frac{1}{2}\right)^x \)；当 \( x < 4 \) 时，\( f(x) = f(x + 1) \). 则 \( f(2 + \log_2 3) = \)（\quad）

\choices
    {\(\frac{1}{24}\)}
    {\(\frac{1}{12}\)}
    {\(\frac{1}{8}\)}
    {\(\frac{3}{8}\)}
\end{problem}

\begin{answer}
A
\end{answer}

\begin{problem}
已知圆 \(C\) 与直线 \(x - y = 0\) 及 \(x - y - 4 = 0\) 都相切，圆心在直线 \(x + y = 0\) 上，则圆 \(C\) 的方程为（\quad）

\choices
    {\((x + 1)^2 +(y - 1)^2 = 2\)}
    {\((x - 1)^{2} + (y + 1)^{2} = 2\)}
    {\((x - 1)^{2} + (y - 1)^{2} = 2\)}
    {\((x + 1)^{2} + (y + 1)^{2} = 2\)}
\end{problem}

\begin{answer}
    B
\end{answer}

\begin{problem}
已知 \(\tan \theta = 2\)，则 \(\sin^2\theta + \sin \theta \cos \theta - 2\cos^2\theta =\)（\quad）

\choices
    {\(-\frac{4}{3}\)}
    {\( \frac{5}{4} \)}
    {\(-\frac{3}{4}\)}
    {\(\frac{4}{5}\)}
\end{problem}

\begin{answer}
    D
\end{answer}

\begin{problem}
\(ABCD\) 为长方形，\(AB = 2\)，\(BC = 1\)，\(O\) 为 \(AB\) 的中点. 在长方形 \(ABCD\) 内随机取一点，取到的点到 \(O\) 的距离大于 \(1\) 的概率为（\quad）

\choices
    {\( \frac{3}{4} \)}
    {\(1 - \frac{\pi}{4}\)}
    {\(\frac{\pi}{8}\)}
    {\(1 - \frac{\pi}{8}\)}
\end{problem}

\begin{answer}
    B
\end{answer}

\begin{problem}
某店一个月的收入和支出总共记录了 \(N\) 个数据 \(a_1\)，\(a_2\)，\(\dots\)，\(a_N\)，其中收入记为正数. 支出记为负数. 该店用下边的程序框图计算月总收入 \(S\) 和月净盈利 \(V\)，那么在图中空白的判断框和处理框中，应分别填入下列四个选项中的（\quad）

\bitmapfigure[max width=0.46\linewidth]{img/2009/liaoning_liberal/q10_fig1.png}

\choices
    {\(A > 0\)，\(V = S - T\)}
    {\(A < 0\)，\(V = S - T\)}
    {\(A > 0\)，\(V = S + T\)}
    {\(A < 0\)，\(V = S + T\)}
\end{problem}

\begin{answer}
    C
\end{answer}

\begin{problem}
下列 \(4\) 个命题：

\begin{circlelist}
\item \(p_1\)：\(\exists x \in (0, +\infty)\)，\(\left(\frac{1}{2}\right)^x < \left(\frac{1}{3}\right)^x\)；
\item \(p_2\)：\(\exists x \in (0,1)\)，\(\log_{\frac12} x > \log_{\frac13} x\)；
\item \(p_3\)：\(\forall x \in (0, +\infty)\)，\(\left(\frac{1}{2}\right)^x > \log_{\frac12} x\)；
\item \(p_4\)：\(\forall x \in \left(0,\frac{1}{3}\right)\)，\(\left(\frac{1}{2}\right)^x < \log_{\frac13} x\).
\end{circlelist}
其中的真命题是（\quad）

\choices
    {\(p_1\)，\(p_3\)}
    {\(p_1\)，\(p_4\)}
    {\(p_2\)，\(p_3\)}
    {\(p_2\)，\(p_4\)}
\end{problem}

\begin{answer}
    D
\end{answer}

\begin{problem}
已知偶函数 \( f(x) \) 在区间 \([0, +\infty)\) 单调增加，则满足 \( f(2x - 1) < f\left(\frac{1}{3}\right) \) 的 \( x \) 取值范围是（\quad）

\choices
    {\(\left(\frac{1}{3},\frac{2}{3}\right)\)}
    {\(\left[\frac{1}{3},\frac{2}{3}\right)\)}
    {\(\left(\frac{1}{2},\frac{2}{3}\right)\)}
    {\(\left[\frac{1}{2},\frac{2}{3}\right)\)}
\end{problem}

\begin{answer}
    A
\end{answer}

\section{填空题}

\begin{problem}
在平面直角坐标系 \(xOy\) 中，四边形 \(ABCD\) 的边 \(AB\parallel DC\)，\(AD\parallel BC\). 已知点 \(A(-2,0)\)，\(B(6,8)\)，\(C(8,6)\)，则 \(D\) 点的坐标为\fillinblank{}.
\end{problem}

\begin{answer}
    \((0,-2)\)
\end{answer}

\begin{problem}
已知函数 \( f(x) = \sin (\omega x + \varphi) \) （\(\omega > 0\)）的图象如图所示，则 \(\omega = \)\fillinblank{}.

\bitmapfigure[max width=0.34\linewidth]{img/2009/liaoning_liberal/q14_fig1.png}
\end{problem}

\begin{answer}
    \(\frac{3}{2}\)
\end{answer}

\begin{problem}
若函数 \( f(x) = \frac{x^2 + a}{x + 1} \) 在 \( x = 1 \) 处取极值，则 \( a = \)\fillinblank{}.
\end{problem}

\begin{answer}
    \(3\)
\end{answer}

\begin{problem}
设某几何体的三视图如下 （尺寸的长度单位为 \(\mathrm{m}\)）. 则该几何体的体积为\fillinblank{}\(\mathrm{m}^{3}\).

\bitmapfigure[max width=0.58\linewidth]{img/2009/liaoning_liberal/q16_fig1.png}
\end{problem}

\begin{answer}
    \(4\)
\end{answer}

\section{解答题}

\begin{problem}
等比数列 \(\{a_{n}\}\) 的前 \(n\) 项和为 \(S_{n}\). 已知 \(S_{1}\)，\(S_{3}\)，\(S_{2}\) 成等差数列.
\begin{enumerate}
\item 求 \(\{a_{n}\}\) 的公比 \(q\)；
\item 若 \(a_{1} - a_{3} = 3\)，求 \(S_{n}\).
\end{enumerate}
\end{problem}

\begin{solution}
\begin{enumerate}
\item
设等比数列的首项为 \(a_{1}\)，公比为 \(q\). 由等比数列的定义，\(a_{1}\neq 0\)，且 \(q\neq 0\). 因为
\(S_{1}=a_{1}\)，\(S_{2}=a_{1}(1+q)\)，\(S_{3}=a_{1}(1+q+q^{2})\).
又 \(S_{1}\)，\(S_{3}\)，\(S_{2}\) 成等差数列，所以
\[
S_{1}+S_{2}=2S_{3}.
\]
代入上式并约去 \(a_{1}\)，得
\[
1+(1+q)=2(1+q+q^{2})
\quad\Longleftrightarrow\quad
q(2q+1)=0.
\]
由于 \(q\neq0\)，故 \(q=-\frac{1}{2}\).

\item
由 \(a_{1}-a_{3}=3\) 及 \(a_{3}=a_{1}q^{2}\)，得
\[
a_{1}\left(1-q^{2}\right)=3
\quad\Longleftrightarrow\quad
a_{1}\left(1-\frac14\right)=3,
\]
所以 \(a_{1}=4\). 因此
\[
S_{n}=\frac{a_{1}(1-q^{n})}{1-q}
    =\frac{4\left(1-\left(-\frac12\right)^{n}\right)}{1+\frac12}
    =\frac{8}{3}\left(1-\left(-\frac12\right)^{n}\right).
\]
\end{enumerate}
\end{solution}

\begin{problem}
如图，\(A\)，\(B\)，\(C\)，\(D\) 都在同一个与水平面垂直的平面内，\(B\)，\(D\) 为两岛上的两座灯塔的塔顶. 测量船于水面 \(A\) 处测得 \(B\) 点和 \(D\) 点的仰角分别为 \(75^{\circ}\)，\(30^{\circ}\)，于水面 \(C\) 处测得 \(B\) 点和 \(D\) 点的仰角均为 \(60^{\circ}\)，\(AC = 0.1 \mathrm{~km}\). 试探究图中 \(B\)，\(D\) 间距离与另外哪两点间距离相等，然后求出 \(B\)，\(D\) 的距离. （计算结果精确到 \(0.01 \mathrm{~km}\)，\(\sqrt{2} \approx 1.414\)，\(\sqrt{6} \approx 2.449\)）

\bitmapfigure[max width=0.42\linewidth]{img/2009/liaoning_liberal/q18_fig1.png}
\end{problem}

\begin{solution}
在 \(\triangle ACD\) 中，\(\angle DAC=30^{\circ}\)，且在 \(C\) 点，射线 \(CA\) 与 \(CD\) 的夹角为 \(180^{\circ}-60^{\circ}=120^{\circ}\)，即 \(\angle ACD=120^{\circ}\). 因而
\[
\angle ADC=180^{\circ}-30^{\circ}-120^{\circ}=30^{\circ}.
\]
所以 \(\angle DAC=\angle ADC\)，从而 \(CD=AC=0.1\mathrm{~km}\)，即 \(\triangle ACD\) 是等腰三角形.

又 \(\angle BCA=60^{\circ}\)，\(\angle BCD=60^{\circ}\)，故 \(CB\) 平分 \(\angle ACD\). 等腰三角形 \(ACD\) 的顶角平分线也是底边 \(AD\) 的垂直平分线，所以
\[
BD=BA.
\]

在 \(\triangle ABC\) 中，\(\angle BAC=180^{\circ}-75^{\circ}=105^{\circ}\)，\(\angle BCA=60^{\circ}\)，\(\angle ABC=15^{\circ}\).
由正弦定理，
\[
\frac{AB}{\sin60^{\circ}}=\frac{AC}{\sin15^{\circ}},
\]
所以
\[
BD=BA=\frac{0.1\sin60^{\circ}}{\sin15^{\circ}}
    =\frac{3\sqrt{2}+\sqrt{6}}{20}\mathrm{~km}
    \approx\frac{3\times1.414+2.449}{20}\mathrm{~km}
    \approx0.33\mathrm{~km}.
\]
\end{solution}

\begin{problem}
如图，已知两个正方形 \(ABCD\) 和 \(DCEF\) 不在同一平面内，\(M\)，\(N\) 分别为 \(AB\)，\(DF\) 的中点.
\begin{enumerate}
\item 若 \(CD = 2\)，平面 \(ABCD \perp\) 平面 \(DCEF\)，求直线 \(MN\) 的长；
\item 用反证法证明：直线 \(ME\) 与 \(BN\) 是两条异面直线.
\end{enumerate}

\bitmapfigure[max width=0.46\linewidth]{img/2009/liaoning_liberal/q19_fig1.png}
\end{problem}

\begin{solution}
\begin{enumerate}
\item 取 \(G\) 为 \(CD\) 的中点，连接 \(MG\)，\(NG\). 由于 \(ABCD\) 为边长为 \(2\) 的正方形，\(MG\parallel AD\)，所以 \(MG\perp CD\)，且 \(MG=AD=2\). 在 \(\triangle DCF\) 中，\(G\)，\(N\) 分别为 \(DC\)，\(DF\) 的中点，故 \(GN\parallel CF\)，从而
    \[
    GN=\frac12CF=\frac12\cdot2=1.
    \]
    因为平面 \(ABCD\perp\) 平面 \(DCEF\)，且 \(MG\) 在平面 \(ABCD\) 内并垂直于两平面的交线 \(CD\)，所以 \(MG\perp\) 平面 \(DCEF\)，进而 \(MG\perp GN\). 因此在直角三角形 \(MGN\) 中，
    \[
    MN=\sqrt{MG^{2}+GN^{2}}=\sqrt{2^{2}+1^{2}}=\sqrt{5}.
    \]
\item 反设直线 \(ME\) 与 \(BN\) 共面，设它们所在的平面为 \(\pi\). 由于 \(M\)，\(B\) 在直线 \(AB\) 上，故 \(AB\subset\pi\). 又 \(E\)，\(N\) 在平面 \(DCEF\) 内，且 \(\pi\) 不等于平面 \(DCEF\)，所以
    \[
    \pi\cap\text{平面 }DCEF=EN.
    \]
    在正方形 \(ABCD\) 中，\(AB\parallel CD\)，而 \(CD\) 是两平面的交线，所以 \(AB\parallel\) 平面 \(DCEF\). 因为 \(AB\subset\pi\)，可得 \(AB\parallel EN\). 又在正方形 \(DCEF\) 中，\(CD\parallel EF\)，故
    \[
    AB\parallel CD\parallel EF.
    \]
    于是 \(EN\parallel EF\). 但直线 \(EN\) 与 \(EF\) 都经过点 \(E\)，这将导致 \(N\in EF\). 同时 \(N\in DF\)，而 \(DF\cap EF=\{F\}\)，故只能有 \(N=F\)，与 \(N\) 是 \(DF\) 的中点矛盾.

    所以直线 \(ME\) 与 \(BN\) 不共面，故它们是异面直线.
\end{enumerate}
\end{solution}

\begin{problem}
某企业有两个分厂生产某种零件，按规定内径尺寸 （单位：\(\mathrm{mm}\)） 的值落在 \([29.94, 30.06]\) 的零件为优质品. 从两个分厂生产的零件中各抽出 \(500\) 件，量其内径尺寸，得结果如下表：

\begin{center}
\begin{tabular}{c@{\hspace{2em}}c}
甲厂 & 乙厂 \\
\begin{tabular}{|c|c|}
\hline
分组 & 频数 \\
\hline
\([29.86,29.90)\) & \(12\) \\
\([29.90,29.94)\) & \(63\) \\
        \([29.94,29.98)\) & \(86\) \\
\([29.98,30.02)\) & \(182\) \\
\([30.02,30.06)\) & \(92\) \\
\([30.06,30.10)\) & \(61\) \\
\([30.10,30.14)\) & \(4\) \\
\hline
\end{tabular}
&
\begin{tabular}{|c|c|}
\hline
分组 & 频数 \\
\hline
\([29.86,29.90)\) & \(29\) \\
\([29.90,29.94)\) & \(71\) \\
        \([29.94,29.98)\) & \(85\) \\
        \([29.98,30.02)\) & \(159\) \\
        \([30.02,30.06)\) & \(76\) \\
        \([30.06,30.10)\) & \(62\) \\
        \([30.10,30.14)\) & \(18\) \\
\hline
\end{tabular}
\end{tabular}
\end{center}

\begin{enumerate}
\item 试分别估计两个分厂生产的零件的优质品率；
\item 由于以上统计数据填下面 \(2 \times 2\) 列联表，并问是否有 \(99\%\) 的把握认为 “两个分厂生产的零件的质量有差异”.
\end{enumerate}
\begin{center}
\begin{tabular}{|c|c|c|c|}
\hline
 & 甲厂 & 乙厂 & 合计 \\
\hline
优质品 &  &  &  \\
\hline
非优质品 &  &  &  \\
\hline
合计 &  &  &  \\
\hline
\end{tabular}
\end{center}

附：\(\chi^2 = \frac{n(n_{11}n_{22} - n_{12}n_{21})^2}{n_{1+}n_{2+}n_{+1}n_{+2}}\)，
\begin{center}
\begin{tabular}{c|cc}
\(P(\chi^2 \geqslant k)\) & \(0.05\) & \(0.01\) \\
\hline
\(k\) & \(3.841\) & \(6.635\)
\end{tabular}
\end{center}
\end{problem}

\begin{solution}
\begin{enumerate}
\item 按表中各组连续分组的题意，优质品对应第三、四、五组. 因此甲厂抽查的优质品有
\[
86+182+92=360\text{ 件},
\]
甲厂生产的零件优质品率估计为
\[
\frac{360}{500}=72\%.
\]
乙厂抽查的优质品有
\[
85+159+76=320\text{ 件},
\]
乙厂生产的零件优质品率估计为
\[
\frac{320}{500}=64\%.
\]
\item 由上面的频数可得列联表为
\begin{center}
\begin{tabular}{|c|c|c|c|}
\hline
 & 甲厂 & 乙厂 & 合计 \\
\hline
优质品 & 360 & 320 & 680 \\
\hline
非优质品 & 140 & 180 & 320 \\
\hline
合计 & 500 & 500 & 1000 \\
\hline
\end{tabular}
\end{center}
所以
\[
\chi^{2}=\frac{1000(360\cdot180-320\cdot140)^{2}}{500\cdot500\cdot680\cdot320}
\approx7.35.
\]
由于 \(7.35>6.635\)，故有 \(99\%\) 的把握认为两个分厂生产的零件的质量有差异.
\end{enumerate}
\end{solution}

\begin{problem}
设 \( f(x) = \e^{x}(ax^{2} + x + 1) \)，且曲线 \( y = f(x) \) 在 \( x = 1 \) 处的切线与 \( x \) 轴平行.
\begin{enumerate}
\item 求 \(a\) 的值，并讨论 \( f(x) \) 的单调性；
\item 证明：当 \(\theta \in \left[0, \frac{\pi}{2}\right]\) 时，\(|f(\cos \theta) - f(\sin \theta)| < 2\).
\end{enumerate}
\end{problem}

\begin{solution}
\begin{enumerate}
\item 由
\[
f'(x)=\e^{x}\left(ax^{2}+x+1+2ax+1\right)
    =\e^{x}\left(ax^{2}+(2a+1)x+2\right),
\]
且曲线在 \(x=1\) 处的切线与 \(x\) 轴平行，得 \(f'(1)=0\). 因为 \(\e\neq0\)，所以
\[
a+(2a+1)+2=0,
\]
即 \(a=-1\). 此时
\[
f'(x)=\e^{x}(-x^{2}-x+2)=-\e^{x}(x+2)(x-1).
\]
由于 \(\e^{x}>0\)，故当 \(x\in(-\infty,-2)\cup(1,+\infty)\) 时，\(f'(x)<0\)；当 \(x\in(-2,1)\) 时，\(f'(x)>0\). 所以 \(f(x)\) 在 \((-\infty,-2)\)，\((1,+\infty)\) 上单调减少，在 \((-2,1)\) 上单调增加.
\item 由上问知 \(f(x)\) 在 \([0,1]\) 上单调增加，且 \(f(0)=1\)，\(f(1)=\e\).
因此对于任意 \(x_{1},x_{2}\in[0,1]\)，有
\[
\left|f(x_{1})-f(x_{2})\right|\leqslant f(1)-f(0)=\e-1<2.
\]
其中 \(\e<3\)，例如由
\[
\e=\sum_{k=0}^{\infty}\frac{1}{k!}<1+1+\sum_{k=2}^{\infty}\frac{1}{2^{k-1}}=3
\]
可得 \(\e-1<2\).
当 \(\theta\in\left[0,\frac{\pi}{2}\right]\) 时，\(\cos\theta,\sin\theta\in[0,1]\)，代入上式即得
\[
\left|f(\cos\theta)-f(\sin\theta)\right|<2.
\]
\end{enumerate}
\end{solution}

\begin{problem}
已知椭圆 \(C\) 过点 \(A\left(1, \frac{3}{2}\right)\)，两个焦点为 \((-1, 0)\)，\((1, 0)\).
\begin{enumerate}
\item 求椭圆 \(C\) 的方程；
\item \(E\)，\(F\) 是椭圆 \(C\) 上的两个动点，如果直线 \(AE\) 的斜率与 \(AF\) 的斜率互为相反数，证明直线 \(EF\) 的斜率为定值. 并求出这个定值.
\end{enumerate}
\end{problem}

\begin{solution}
\begin{enumerate}
\item 设椭圆的长半轴、短半轴分别为 \(a\)，\(b\)，焦距的一半为 \(c\). 由焦点坐标可知 \(c=1\)，所以
\[
a^{2}-b^{2}=c^{2}=1.
\]
又因为点 \(A\left(1,\frac32\right)\) 在椭圆上，故
\[
\frac{1}{a^{2}}+\frac{9}{4b^{2}}=1.
\]
代入 \(b^{2}=a^{2}-1\)，得
\[
\frac{1}{a^{2}}+\frac{9}{4(a^{2}-1)}=1
\quad\Longleftrightarrow\quad
4a^{4}-17a^{2}+4=0.
\]
于是
\[
(a^{2}-4)(4a^{2}-1)=0.
\]
因为 \(a^{2}>1\)，所以 \(a^{2}=4\)，进而 \(b^{2}=3\). 因此椭圆 \(C\) 的方程为
\[
\frac{x^{2}}{4}+\frac{y^{2}}{3}=1.
\]
\item 设直线 \(AE\) 的斜率为 \(k\)，则其方程为
\[
y=k(x-1)+\frac32.
\]
代入椭圆方程并乘以 \(12\)，得
\[
(3+4k^{2})x^{2}+4k(3-2k)x+4\left(\frac32-k\right)^{2}-12=0.
\]
由于点 \(A\) 在该直线上且横坐标为 \(1\)，上式的一个根为 \(1\). 设另一个交点为 \(E(x_{E},y_{E})\)，由根与系数的关系得
\[
x_{E}=\frac{4\left(\frac32-k\right)^{2}-12}{3+4k^{2}}.
\]
又 \(y_{E}=kx_{E}+\frac32-k\).
直线 \(AF\) 的斜率为 \(-k\)，同理
\[
x_{F}=\frac{4\left(\frac32+k\right)^{2}-12}{3+4k^{2}}.
\]
又 \(y_{F}=-kx_{F}+\frac32+k\).
由于直线 \(EF\) 有定义，\(k\neq0\). 直接计算得
\[
x_{F}-x_{E}=\frac{24k}{3+4k^{2}}.
\]
以及
\[
2-x_{E}-x_{F}=\frac{12}{3+4k^{2}}.
\]
所以
\[
\frac{y_{F}-y_{E}}{x_{F}-x_{E}}
=\frac{-k(x_{F}+x_{E})+2k}{x_{F}-x_{E}}
=\frac{k(2-x_{E}-x_{F})}{x_{F}-x_{E}}
=\frac12.
\]
因此直线 \(EF\) 的斜率为定值 \(\frac12\).
\end{enumerate}
\end{solution}
